\documentclass[11pt]{amsart}
\usepackage[margin=1.05in]{geometry}
\usepackage{amsmath,amssymb,amsthm}
\usepackage{booktabs}
\usepackage{enumitem}
\usepackage[hidelinks]{hyperref}

\newtheorem{theorem}{Theorem}[section]
\newtheorem{lemma}[theorem]{Lemma}
\newtheorem{proposition}[theorem]{Proposition}
\newtheorem{corollary}[theorem]{Corollary}
\newtheorem*{theoremA}{Theorem A}
\newtheorem*{theoremB}{Theorem B}
\newtheorem*{theoremC}{Theorem C}
\newtheorem*{theoremD}{Theorem D}
\theoremstyle{definition}
\newtheorem{definition}[theorem]{Definition}
\newtheorem{question}{Question}
\theoremstyle{remark}
\newtheorem{remark}[theorem]{Remark}

\newcommand{\Z}{\mathbb{Z}}
\newcommand{\N}{\mathbb{N}}
\newcommand{\pos}{\operatorname{pos}}

\title[Symmetric insertion and 2-adic completions on $\Z$]{Symmetric insertion and 2-adic completions for permutations of $\Z$:\\ odd differences, and obstructions towards Erd\H{o}s Problem 195}
\author{Keio}
\date{October 29, 2026}

\begin{document}

\begin{abstract}
A permutation of $\Z$ is a bijection $\N\to\Z$, written as a sequence $a_0,a_1,\dots$; a monotone $k$-term arithmetic progression (monotone $k$-AP) is a subsequence $a_{i_1},\dots,a_{i_k}$, $i_1<\dots<i_k$, forming an increasing or decreasing arithmetic progression. Erd\H{o}s Problem 195 asks for the largest $k$ such that every permutation of $\Z$ contains a monotone $k$-AP; it is known that $k\in\{3,4\}$. After Ho's recent construction of a permutation of $\N$ without monotone $4$-APs by means of $2$-adic completions (``safe words''), the case of $\Z$ is the natural next step. This paper does \emph{not} settle Problem 195. We prove three results that delimit what $2$-adic methods can do on $\Z$.

(i) \emph{Symmetric insertion.} There is an explicit permutation of $\Z$ in which every even integer is preceded by a set of odd integers that is symmetric about it. Consequently it contains no monotone $4$-AP with odd common difference, and every monotone $3$-AP with odd difference has an odd middle term. More generally, for every fixed $k\ge0$ there is a permutation of $\Z$ without monotone $4$-APs whose difference $d$ satisfies $v_2(d)=k$.

(ii) \emph{Two-level obstruction.} Symmetric insertion cannot be imposed (up to finitely many exceptions) both at a node of the $2$-adic tree and at its unprotected child, nor at a node and at the two grandchildren below its unprotected child; so this mechanism alone cannot even exclude all monotone $4$-APs with difference $\not\equiv0\pmod 4$. (Inside the protected child it can be iterated.)

(iii) \emph{Completion obstructions.} No permutation of $\Z$ has all sufficiently long prefixes safe with respect to $2$-adic tree completions. If a safe word has the block form $U\,A\,B$ or $U\,B\,A$, where $A$ has the preferred parity and $B\neq\emptyset$ the other parity, some entry of the preferred parity precedes $B$, and the entries of $U$ of the other parity lie between the extreme entries of the preferred parity preceding $B$, then $B$ consists of one integer, the midpoint of these extremes. As a consequence, the parity-block extension pattern underlying Ho's construction, with a fixed root preference, cannot exhaust $\Z$.
\end{abstract}

\maketitle

\section{Introduction}

Davis, Entringer, Graham and Simmons \cite{DEGS} showed that every permutation of $\N$ or of $\Z$ contains a monotone $3$-AP and that monotone $5$-APs (for $\N$) can be avoided. (For density versions of these questions see \cite{LV11,Ge26}.) For $\Z$, Geneson \cite{Ge19} and Adenwalla \cite{Ad22} constructed permutations avoiding monotone $6$-APs and $5$-APs respectively. Whether monotone $4$-APs can be avoided was open for both $\N$ and $\Z$; for $\N$ it was recently settled by Ho \cite{Ho26}, who constructed a permutation of the positive integers with no monotone $4$-AP. The corresponding question for $\Z$ is Erd\H{o}s Problem~195 \cite{EP195}:

\begin{quote}
\emph{What is the largest $k$ such that in any permutation of $\Z$ there must exist a monotone $k$-term arithmetic progression?}
\end{quote}

Here a permutation of $\Z$ is a one-sided arrangement $a_0,a_1,a_2,\dots$ of all integers (as in the formalisation \cite{FC195}); the answer is $3$ or $4$. A permutation of $\Z$ without monotone $4$-APs restricts to one of $\N$, so the $\Z$-problem is at least as hard as the $\N$-problem.

Ho's method uses the $2$-adic order $\lhd$ (compare binary digits from the least significant end) and \emph{safe words}: finite words $W$ such that $W$ followed by all remaining integers in $\lhd$-order contains no monotone $4$-AP. The construction extends safe words recursively by splitting into parity classes and appending a block of new odd integers followed by a block of new even integers. In $\N$ the decisive point is that a reflection $2y-x$ of a large odd $x$ through a small even $y$ is negative and therefore irrelevant. In $\Z$ such reflections always exist, and this paper studies what this changes.

\subsection*{Results}
For a node of the binary tree of residue classes (a class $r\bmod 2^k$) we single out one child as \emph{protected} and call an element $q$ of the protected child \emph{protected} if the elements of the other child that precede $q$ form a set symmetric about $q$ (Definition~\ref{def:prot}). Protection kills all in-order $3$-APs whose middle term is $q$ and whose outer terms lie in the sibling class (Lemma~\ref{lem:kill}).

\begin{theoremA}[Theorem~\ref{thm:odd} and Corollary~\ref{cor:level}]
There is an explicit permutation of $\Z$ in which every even integer is protected for the root. It contains no monotone $4$-AP with odd common difference, and every monotone $3$-AP with odd difference has an odd middle term. For every $k\ge0$ there is a permutation of $\Z$ without monotone $4$-APs whose difference $d$ has $v_2(d)=k$.
\end{theoremA}

For the positive integers, LeSaulnier and Vijay \cite{LV11} constructed a permutation without monotone $4$-APs of odd difference, and Adenwalla \cite{Ad23} extended this to differences not divisible by $2^k$. We have not found Theorem~A for $\Z$ in the literature; its proof is elementary.

\begin{theoremB}[Theorems~\ref{thm:twolevel} and~\ref{thm:level02}]
Let one parity class be the protected child of the root, and let $P$ be the other parity class. In no permutation of $\Z$ are all but finitely many elements of the protected parity class protected for the root and, at the same time, all but finitely many elements of one residue class modulo $4$ contained in $P$ protected for the node $P$. The same holds with the two level-$2$ nodes contained in $P$ (each with one protected residue class modulo $8$) in place of the node $P$.
\end{theoremB}

Thus the mechanism of Theorem~A cannot be iterated from a node to its unprotected child, nor from a node to the two grandchildren below its unprotected child. It can be iterated into the protected child (Remark~\ref{rem:inside}); but excluding all differences $\not\equiv0\pmod4$ by protection would require protection at both children of the root. Our third group of results concerns $2$-adic completions with arbitrary tree orders (all node-wise choices of a preferred child, Section~\ref{sec:tree}).

\begin{theoremC}[Proposition~\ref{prop:allprefix}]
For every permutation of $\Z$, infinitely many prefixes are not safe for any tree order.
\end{theoremC}

\begin{theoremD}[Theorem~\ref{thm:block} and Corollary~\ref{cor:ho}]
Let $\prec$ be a tree order whose root prefers parity $p$. If $W=U\,A\,B$ or $W=U\,B\,A$ is $\prec$-safe, where $A$ has parity $p$, $B\neq\emptyset$ has parity $1-p$, some entry of parity $p$ precedes $B$, and the parity-$(1-p)$ entries of $U$ lie between the extreme parity-$p$ entries preceding $B$, then $B$ consists of a single integer, the midpoint of those extremes. Consequently no sequence of safe words obtained by repeatedly appending a parity-$p$ block and then a parity-$(1-p)$ block, with a fixed root preference $p$, exhausts $\Z$.
\end{theoremD}

Ho's construction needs safety only at the checkpoints of the recursion, so Theorem~C alone does not rule it out; Theorem~D addresses the checkpoints of the parity-block scheme directly. Neither theorem excludes every conceivable $2$-adic construction for $\Z$. Lemma~\ref{lem:largest} records the constraint that any successful scheme must respect. Problem~195 remains open.

\section{Preliminaries}\label{sec:prelim}

\subsection{Sequences and in-order progressions}
Let $\sigma=(a_t)_{t\ge0}$ be a sequence of distinct integers, and write $\pos(x)=t$ if $a_t=x$. A triple $(x,y,z)$ with $y-x=z-y\ne0$ is \emph{in order} if all three terms occur and $\pos(x)<\pos(y)<\pos(z)$. A $k$-AP is \emph{monotone} if its terms, listed in one of its two directions, occur at increasing positions; listed this way we call it \emph{in order}. Every in-order $4$-AP $(a,b,c,d)$ contains the in-order $3$-APs $(a,b,c)$ and $(b,c,d)$.

\subsection{Residue classes}
For $k\ge0$ and $r\in\Z$ let $N=(k,r)$ denote the class $r+2^k\Z$ (a \emph{node} of level $k$). Its two \emph{children} are the classes $r+2^{k+1}\Z$ and $r+2^k+2^{k+1}\Z$. If $d\ne0$ and $v_2(d)=k$, the terms of any progression with difference $d$ lie in a single node of level $k$ and alternate between its two children. The affine bijection $x\mapsto (x-r)/2^k$ maps the node $(k,r)$ onto $\Z$, maps progressions to progressions, and maps the two children to the even and the odd integers.

\section{Symmetric insertion}\label{sec:sym}

\begin{definition}\label{def:prot}
Let $N$ be a node and choose one of its children, $Q_N$, as the \emph{protected child}; call the other child $P_N$. In a sequence $\sigma$, an element $q\in Q_N$ is \emph{$N$-protected} if the set
\[
S_N(q)=\{x\in P_N:\ x \text{ occurs in }\sigma\text{ and }\pos(x)<\pos(q)\}
\]
is symmetric about $q$, i.e.\ $2q-S_N(q)=S_N(q)$.
\end{definition}

\begin{lemma}\label{lem:kill}
If $q\in Q_N$ is $N$-protected, then there is no in-order $3$-AP $(x,q,z)$ with $x,z\in P_N$.
\end{lemma}

\begin{proof}
If $x\in P_N$ precedes $q$, then $x\in S_N(q)$, so $z=2q-x\in S_N(q)$ precedes $q$.
\end{proof}

\begin{lemma}\label{lem:levelk}
Let $N$ be a node of level $k$ and suppose every element of $Q_N$ occurring in $\sigma$ is $N$-protected. Then $\sigma$ contains no monotone $4$-AP with terms in $N$ and difference $d$ satisfying $v_2(d)=k$.
\end{lemma}

\begin{proof}
List such a progression in order as $(x_1,x_2,x_3,x_4)$. Its terms alternate between the children of $N$, so either $x_2\in Q_N$ and $x_1,x_3\in P_N$, or $x_3\in Q_N$ and $x_2,x_4\in P_N$. The in-order $3$-AP $(x_1,x_2,x_3)$, respectively $(x_2,x_3,x_4)$, contradicts Lemma~\ref{lem:kill}.
\end{proof}

\begin{theorem}\label{thm:odd}
Let $y_1,y_2,\dots$ be an enumeration of the even integers and $t_1,t_2,\dots$ an enumeration of the odd integers. Build words $W_0=\emptyset\subset W_1\subset\cdots$, each a prefix of the next, as follows. Given $W_{j-1}$, first append $t_j$ unless it already occurs. Let $O$ be the set of odd integers occurring at that moment. Then append the integers of $(2y_j-O)\setminus O$ in increasing order, and then append $y_j$. The limit sequence $\sigma$ is a permutation of $\Z$ in which every even integer is protected for the root node (with $Q=$ even integers, $P=$ odd integers). Consequently:
\begin{enumerate}[label=\textup{(\alph*)}]
\item $\sigma$ has no in-order $3$-AP of the form (odd, even, odd); in particular every monotone $3$-AP of odd difference has an odd middle term;
\item $\sigma$ has no monotone $4$-AP with odd common difference.
\end{enumerate}
\end{theorem}

\begin{proof}
Even integers are appended only as the targets $y_j$, at the end of step $j$; the reflected integers $2y_j-O$ and the $t_j$ are odd. Hence no integer is appended twice. Each step appends finitely many integers, and $t_j$ and $y_j$ occur in $W_j$; so $\sigma$ is a permutation of $\Z$. At that moment the set $O'$ of odd integers already present is $O\cup(2y_j-O)$, where $O$ is the set before the reflections were appended, so $2y_j-O'=(2y_j-O)\cup O=O'$. The set $O'$ is exactly $S_N(y_j)$ for the root node $N$, so $y_j$ is protected. (a) is Lemma~\ref{lem:kill}, and (b) is Lemma~\ref{lem:levelk} with $k=0$.
\end{proof}

\begin{remark}[The first terms]
With $y=0,2,-2,4,-4,6,\dots$ and $t=1,-1,3,-3,\dots$ the sequence begins
\[
\begin{aligned}
&1,-1,\mathbf{0},3,5,\mathbf{2},-9,-7,-5,-3,\mathbf{-2},7,9,11,13,15,17,\mathbf{4},\\
&-25,-23,\dots,-11,\mathbf{-4},19,21,\dots,37,\mathbf{6},\dots
\end{aligned}
\]
(even terms in bold). For instance, the odd integers preceding $-2$ are $\{1,-1,3,5,-9,-7,-5,-3\}$, which is symmetric about $-2$. For these enumerations one checks by induction that after step $n$ the odd integers present are exactly the odd integers of an interval symmetric about $y_n$, that step $n\ge2$ adds $|y_n-y_{n-1}|=2(n-1)$ odd integers and $y_n$, and hence that $W_n$ has $n^2+2$ terms (we also confirmed this by computer for $n\le300$).
\end{remark}

\begin{corollary}\label{cor:level}
For each $k\ge0$ there is a permutation of $\Z$ with no monotone $4$-AP whose difference $d$ satisfies $v_2(d)=k$.
\end{corollary}

\begin{proof}
Let $\tau$ be the permutation of Theorem~\ref{thm:odd}. For each residue $r\in\{0,\dots,2^k-1\}$ let $\tau_r$ be the image of $\tau$ under $x\mapsto r+2^kx$; it lists the class $r+2^k\Z$. Interleave $\tau_0,\dots,\tau_{2^k-1}$ in round-robin fashion. A progression with $v_2(d)=k$ lies in one class $r+2^k\Z$, interleaving preserves relative positions within the class, and $x\mapsto(x-r)/2^k$ turns it into a progression of odd difference in $\tau$.
\end{proof}

\begin{remark}
Theorem~\ref{thm:odd} only controls progressions of odd difference: in the sequence of Theorem~\ref{thm:odd} monotone $4$-APs of even difference are abundant (Section~\ref{sec:comp}). The construction is a $\Z$-analogue of the odd-difference results for $\N$ in \cite{LV11,Ad23}. On $\Z$ the reflections $2y-O$ always exist, so the symmetric sets must be built explicitly; the reflected sets grow only linearly in this construction (Section~\ref{sec:comp}).
\end{remark}

\section{A node and its unprotected child cannot both be protected}\label{sec:twolevel}

The natural attempt to treat all differences is to protect one child at every node. The following theorem shows that this already fails at levels $0$ and $1$.

\begin{theorem}\label{thm:twolevel}
Let $\sigma$ be a permutation of $\Z$. At the root choose the protected child $Q_0$ (one parity class) and let $P_0$ be the other parity class. At the level-$1$ node $P_0$ choose the protected child $Q_1$ (one residue class modulo $4$ inside $P_0$) and let $P_1=P_0\setminus Q_1$. Then it is impossible that all but finitely many elements of $Q_0$ are protected for the root and all but finitely many elements of $Q_1$ are protected for the node $P_0$.
\end{theorem}

\begin{proof}
Suppose both hold, and let $F$ be the finite set of exceptions. Choose $T$ so large that $a_0,\dots,a_{T}$ contain $F$ and at least one element of $P_1$. Put $m_t=\min\{a_0,\dots,a_t\}$ and $M_t=\max\{a_0,\dots,a_t\}$. Since $\sigma$ is onto $\Z$, there is $T'\ge T$ such that for $t\ge T'$ neither $m_t$ nor $M_t$ is among $a_0,\dots,a_T$.

\emph{Claim 1. For $t\ge T'$ we have $m_t,M_t\in P_1$.} Let $e=M_t$; it occurs after $a_T$, so $e\notin F$, and every element preceding $e$ is smaller than $e$. If $e\in Q_0$, then $e$ is protected, and $S_{\rm root}(e)$ is a nonempty set (it contains the elements of $P_1\subseteq P_0$ among $a_0,\dots,a_T$) of integers smaller than $e$ that is symmetric about $e$, which is impossible. Hence $e\in P_0$. If $e\in Q_1$, the same argument applied to $S_{P_0}(e)$, which contains an element of $P_1$ among $a_0,\dots,a_T$, gives a contradiction. So $e\in P_1$. The argument for $m_t$ is symmetric.

\emph{Claim 2. No element of $Q_0$ occurring at a position $t+1>T'$ is protected.} Let $y=a_{t+1}\in Q_0$ be protected, and let $S=S_{\rm root}(y)$, the set of elements of $P_0$ among $a_0,\dots,a_t$. By Claim~1, $m_t,M_t\in P_1\subseteq P_0$, so $\min S=m_t$ and $\max S=M_t$. Symmetry of $S$ about $y$ gives $m_t+M_t=2y$. Since $m_t,M_t$ lie in the same class $P_1$ modulo $4$, $m_t+M_t\equiv 2m_t\pmod 4$, hence $y=(m_t+M_t)/2\equiv m_t\pmod 2$, i.e.\ $y\in P_0$, a contradiction.

Infinitely many elements of $Q_0$ occur at positions $>T'+1$, and all but finitely many of them are protected; this contradicts Claim~2.
\end{proof}

\begin{remark}
Applying the affine renormalisation of a node, the same conclusion holds for any node $N$ of level $k$ and the child $N'=P_N$ of $N$ that is \emph{not} protected: in a permutation of $\Z$, protection cannot hold, up to finitely many exceptions, both at $N$ and at $N'$. To exclude all monotone $4$-APs with difference $\not\equiv0\pmod 4$ by symmetric insertion one would need protection at the root and at both nodes of level $1$, in particular at the unprotected child of the root. So symmetric insertion alone cannot produce such a permutation, and a fortiori cannot settle Problem~195.
\end{remark}

\begin{remark}\label{rem:inside}
The protected child behaves differently. Let $\tau=(\tau_j)_{j\ge1}$ be the sequence of Theorem~\ref{thm:odd}, and run the construction of Theorem~\ref{thm:odd} again with the even integers enumerated as $y_j=2\tau_j$. In the resulting permutation $\sigma_1$ every even integer is protected for the root, and the even integers occur in the order $2\tau_1,2\tau_2,\dots$; so the elements $\equiv2\pmod 4$ preceding an element $2\tau_j\equiv0\pmod4$ form the set $2\,S(\tau_j)$, where $S(\tau_j)$ is the (symmetric) set of odd integers preceding $\tau_j$ in $\tau$. Hence every element of $4\Z$ is also protected for the node $2\Z$, and by Lemma~\ref{lem:levelk} $\sigma_1$ has no monotone $4$-AP of odd difference and none with difference $\equiv2\pmod4$ inside $2\Z$. Monotone $4$-APs with difference $\equiv2\pmod4$ inside the odd integers are not controlled.
\end{remark}

The obstruction is not limited to consecutive levels.

\begin{theorem}\label{thm:level02}
Let $\sigma$ be a permutation of $\Z$. At the root choose a protected child $Q_0$, let $P_0$ be the other parity class, and at each of the two level-$2$ nodes contained in $P_0$ (the two residue classes modulo $4$ inside $P_0$) choose a protected child (a residue class modulo $8$). Then it is impossible that all but finitely many elements of $Q_0$ are protected for the root and all but finitely many elements of the two chosen classes modulo $8$ are protected for their level-$2$ nodes.
\end{theorem}

\begin{proof}
For a level-$2$ node $C\subseteq P_0$ let $\pi_C$ denote its unprotected child, a residue class modulo $8$. Suppose the conclusion fails, let $F$ be the finite set of exceptions, and choose $T$ such that $a_0,\dots,a_T$ contain $F$ and at least one element of each of the four residue classes modulo $8$ contained in $P_0$. Define $m_t,M_t$ and $T'$ as in the proof of Theorem~\ref{thm:twolevel}. As there (Claim~1), for $t\ge T'$ the extremes $m_t,M_t$ lie in $P_0$: a protected element of $Q_0$ cannot be an extreme, since the elements of $P_0$ preceding it form a nonempty symmetric set. Moreover each of $m_t,M_t$ lies in the unprotected child of its level-$2$ node, by the same argument applied to that node. Now let $y\in Q_0$ occur at a position $t+1>T'$ and be protected. As in Claim~2, $y=(m_t+M_t)/2$. Since $y$ and the elements of $P_0$ have different parities, $m_t+M_t=2y\equiv 2m_t+2\pmod 4$, so $m_t\not\equiv M_t\pmod 4$. Hence $m_t$ and $M_t$ lie in the two different level-$2$ nodes $C,C'$ contained in $P_0$, so $m_t\in\pi_C$ and $M_t\in\pi_{C'}$ (or vice versa). Then $m_t+M_t\pmod 8$ is the fixed residue $\pi_C+\pi_{C'}$, and $y$ lies in one fixed residue class modulo $4$. But $Q_0$ contains two residue classes modulo $4$, each with infinitely many elements occurring after position $T'+1$, and all but finitely many of them are protected. Contradiction.
\end{proof}

\section{Tree orders and safe words}\label{sec:tree}

A \emph{tree order} $\prec$ on $\Z$ is given by choosing, at every node $(k,r)$, one of its two children as the \emph{preferred} one; for $x\ne y$ with $k=v_2(x-y)$, we set $x\prec y$ if $x$ lies in the preferred child of the node $(k,x\bmod 2^k)$. This is a linear order (lexicographic order of $2$-adic digits with node-dependent digit orders). Ho's order $\lhd$ \cite{Ho26} is the tree order preferring the child with digit $1$ at every node. We say that $\prec$ \emph{prefers parity $p$ at the root} if the class $p+2\Z$ is the preferred child of the root. The following two facts are standard \cite{ABJ,Ho26}.

\begin{lemma}\label{lem:tree}
Let $\prec$ be a tree order.
\begin{enumerate}[label=\textup{(T\arabic*)}]
\item No $3$-AP is $\prec$-monotone.
\item If $(a,b,c,d)$ is a $4$-AP, then $a\prec b$ if and only if $c\prec d$.
\end{enumerate}
\end{lemma}

\begin{proof}
Let the difference be $\delta$, $k=v_2(\delta)$. All terms lie in one node of level $k$; consecutive terms lie in different children and terms two apart in the same child. So every comparison of consecutive terms is decided at that node by the preference there. For (T1), the middle term of $(a,b,c)$ is in the other child than $a$ and $c$, hence is the $\prec$-largest or the $\prec$-smallest of the three. For (T2), $a$ and $c$ lie in one child and $b,d$ in the other.
\end{proof}

\begin{definition}
Let $\prec$ be a tree order and $W$ a finite word of distinct integers. The \emph{completion} $C_\prec(W)$ is the linear order on $\Z$ that lists $W$ first and then $\Z\setminus W$ in $\prec$-order. The word $W$ is \emph{$\prec$-safe} if $C_\prec(W)$ contains no monotone $4$-AP, i.e.\ no $4$-AP whose terms, in one of the two directions, increase in $C_\prec(W)$.
\end{definition}

If $W_1\subset W_2\subset\cdots$ are words, each a prefix of the next, each $W_j$ is $\prec_j$-safe for some tree order $\prec_j$, and every integer occurs in some $W_j$, then the limit sequence is a permutation of $\Z$ without monotone $4$-APs: four given integers lie in some $W_j$, and their relative order in the limit equals that in $C_{\prec_j}(W_j)$. This is the method of \cite{Ho26} for $\N$.

\begin{lemma}\label{lem:escape}
Let $W$ be $\prec$-safe, and let $x$ precede $y$ in $W$ with $x\prec y$. Put $z=2y-x$ and $w=3y-2x$. Then $z$ occurs in $W$ before $y$, or $w$ occurs in $W$.
\end{lemma}

\begin{proof}
Otherwise $w\notin W$, and either $z\notin W$ or $z$ occurs after $y$. In the first case $z,w$ lie in the tail and $z\prec w$ by (T2), so $(x,y,z,w)$ is monotone in $C_\prec(W)$. In the second case $(x,y,z)$ occurs in this order in $W$ and $w$ lies in the tail, so again $(x,y,z,w)$ is monotone.
\end{proof}

\begin{lemma}\label{lem:extreme}
Let $W$ be $\prec$-safe and $M=\max W$. Then $M\prec x$ for every $x$ that precedes $M$ in $W$. The same holds for $\min W$.
\end{lemma}

\begin{proof}
If $x$ precedes $M$ and $x\prec M$, Lemma~\ref{lem:escape} requires $2M-x$ or $3M-2x$ to occur in $W$; both exceed $M$. The minimum is treated in the same way.
\end{proof}

\begin{proposition}\label{prop:allprefix}
Let $(a_t)_{t\ge0}$ be a permutation of $\Z$ and $P_t=(a_0,\dots,a_t)$. Then for infinitely many $t$ the word $P_t$ is not $\prec$-safe for any tree order $\prec$.
\end{proposition}

\begin{proof}
Suppose that for every $t\ge t_0$ the word $P_t$ is $\prec_t$-safe for some tree order $\prec_t$; let $p_t$ be the parity preferred by $\prec_t$ at the root. Let $m_t=\min P_t$, $M_t=\max P_t$. Choose $T\ge t_0$ such that $P_T$ meets all four residue classes modulo $4$, and $T'\ge T$ such that for $t\ge T'$ neither $m_t$ nor $M_t$ occurs in $P_T$.

(i) \emph{For $t\ge T'$, $m_t\equiv M_t\equiv p_t\pmod 2$.} Some $x\in P_T$ of parity different from $M_t$ precedes $M_t$; by Lemma~\ref{lem:extreme}, $M_t\prec_t x$, which is decided at the root, so $M_t\equiv p_t$. Likewise for $m_t$.

(ii) \emph{For $t\ge T'$, $m_t\equiv M_t\pmod 4$.} Let $e$ be the one of $m_t,M_t$ that occurs first and $f$ the other. Lemma~\ref{lem:extreme} for $f$ gives $f\prec_t e$. If $e\not\equiv f\pmod 4$, this comparison is decided at the level-$1$ node containing $e$ and $f$, whose preferred child is then the class of $f$ modulo $4$. But $e$ is preceded by some $x\in P_T$ with $x\equiv f\pmod 4$, and Lemma~\ref{lem:extreme} for $e$ gives $e\prec_t x$, decided at the same node, so the class of $e$ would be preferred. Contradiction.

(iii) \emph{For $t>T'$, $a_t\equiv p_t\pmod 2$.} Otherwise $a_t$ is not an extreme of $P_t$ by (i), so $m_t<a_t<M_t$ and both $m_t,M_t$ precede $a_t$. As $m_t\prec_t a_t$ (root), Lemma~\ref{lem:escape} gives $2a_t-m_t\le M_t$ (if $2a_t-m_t$ precedes $a_t$) or $3a_t-2m_t\le M_t$, which implies $2a_t-m_t<M_t$. Hence $2a_t\le m_t+M_t$. Symmetrically, using $M_t\prec_t a_t$, $2a_t\ge m_t+M_t$. So $a_t=(m_t+M_t)/2$, which by (ii) has the parity of $m_t$, i.e.\ $p_t$. Contradiction.

(iv) \emph{For $t\ge T'$, $p_{t+1}=p_t$.} If $a_{t+1}$ is not a new maximum then $M_{t+1}=M_t$, and $p_{t+1}\equiv M_{t+1}=M_t\equiv p_t$ by (i); otherwise $m_{t+1}=m_t$ and the same argument applies with minima.

By (iii) and (iv) all $a_t$ with $t>T'$ have the same parity, which is absurd.
\end{proof}

Ho's construction on $\N$ only requires safety at the checkpoints of its recursion, and the intermediate prefixes need not be safe; so Proposition~\ref{prop:allprefix} does not by itself exclude an adaptation of the method to $\Z$. The next theorem concerns the checkpoints of the parity-block scheme.

\begin{theorem}\label{thm:block}
Let $\prec$ be a tree order preferring parity $p$ at the root. Let $W=U\,A\,B$ or $W=U\,B\,A$ be a concatenation of finite words in which every entry of $A$ has parity $p$ and $B$ is nonempty with every entry of parity $1-p$. Suppose that some entry of parity $p$ occurs before the block $B$ (in $UA$, respectively in $U$), let $L$ and $R$ be the smallest and the largest such entry, and assume that every entry of $U$ of parity $1-p$ lies in $[L,R]$. If $W$ is $\prec$-safe, then $(L+R)/2$ is an integer of parity $1-p$ and $B=\{(L+R)/2\}$.
\end{theorem}

\begin{proof}
Let $y\in B$ and suppose $2y>L+R$. Since $L$ precedes $y$ and $L\prec y$ (decided at the root), Lemma~\ref{lem:escape} shows that $z=2y-L$ precedes $y$ or $w=3y-2L$ occurs in $W$. The integer $z$ has parity $p$ and $z>R$, while every entry of parity $p$ preceding $y$ lies in $[L,R]$; so $z$ does not precede $y$, and $w\in W$. Now $w$ has parity $1-p$ and $w>z>R$, so $w$ is neither in $U$ (by hypothesis) nor in $A$ (by parity); hence $w\in B$ and $w>y$, and $2w>L+R$. Iterating produces an infinite strictly increasing sequence in the finite set $B$, a contradiction. Hence $2y\le L+R$, and the same argument with $R$ in place of $L$ gives $2y\ge L+R$.
\end{proof}

\begin{corollary}\label{cor:ho}
Fix $p\in\{0,1\}$. Let $W_1\subset W_2\subset\cdots$ be finite words with $W_{j+1}=W_jA_jB_j$, where $A_j$ consists of integers of parity $p$ and $B_j$ of integers of parity $1-p$, and suppose each $W_j$ ($j\ge2$) is safe for some tree order preferring parity $p$ at the root. Then $\bigcup_jW_j\ne\Z$.
\end{corollary}

\begin{proof}
Suppose $\bigcup_j W_j=\Z$, and let $(a_t)$ be the limit permutation. Let $s$ be the position of the first entry of parity $p$, and $T>s$ such that $a_0,\dots,a_T$ meet all residue classes modulo $4$.

\emph{Step 1.} For large $j$ the hypotheses of Theorem~\ref{thm:block} hold for $U=W_j$, $A=A_j$, $B=B_j$ whenever $B_j\ne\emptyset$. Indeed, for large $j$ the word $W_j$ contains $a_s$, an entry of parity $p$. By Lemma~\ref{lem:extreme}, the maximum of the safe word $W_j$ is either of parity $p$ or preceded by no entry of parity $p$, i.e.\ among $a_0,\dots,a_{s-1}$; for large $j$ the latter is impossible, since entries of parity $p$ are unbounded. Hence every entry of $W_j$ of parity $1-p$ is at most the largest entry of parity $p$ of $W_j$, and similarly from below.

\emph{Step 2.} For large $j$, Theorem~\ref{thm:block} gives $B_j\subseteq\{(L_j+R_j)/2\}$, with $L_j,R_j$ the extreme entries of parity $p$ of $W_jA_j$ (trivially if $B_j=\emptyset$). By Step~1 these are also the extremes of $W_{j+1}$.

\emph{Step 3.} For large $j$, $L_j\equiv R_j\pmod 4$: this is part (ii) of the proof of Proposition~\ref{prop:allprefix}, applied to the safe word $W_{j+1}$ and its extremes $L_j,R_j$, which for large $j$ occur after $a_T$. Hence $(L_j+R_j)/2$ has parity $p$ and $B_j=\emptyset$ for all large $j$.

Then only finitely many integers of parity $1-p$ occur, contradicting $\bigcup_jW_j=\Z$.
\end{proof}

\begin{remark}
In Ho's construction for $\N$ \cite[Algorithm~1]{Ho26} the extension step has the form $P\cdot O\cdot E$ with $O$ new odd and $E$ new even integers, and $\lhd$ prefers odd integers at the root, so it is of the form considered in Corollary~\ref{cor:ho} at the top level. On $\N$ the argument of Theorem~\ref{thm:block} breaks down at the step that uses the largest preferred element $R$: for $y\in E$ the progression $(R,y,2y-R,3y-2R)$ may leave $\N$, and then the analogue of Lemma~\ref{lem:escape} imposes no condition. On $\Z$ every such progression exists.
\end{remark}

Theorem~\ref{thm:block} uses that the escape $w$ must lie in the block $B$. In general an escape of parity $1-p$ may have been placed earlier, before the elements it serves. The following lemma states the constraint that remains for every safe word.

\begin{lemma}\label{lem:largest}
Let $W$ be safe for a tree order preferring parity $p$ at the root. Let $y$ be the largest entry of $W$ of parity $1-p$, and suppose some entry of parity $p$ precedes $y$; let $L$ and $R$ be the least and largest entries of parity $p$ preceding $y$. Then $2y\le L+R$. Symmetrically, if $y'$ is the least entry of $W$ of parity $1-p$ and some entry of parity $p$ precedes $y'$, then $2y'\ge L'+R'$, where $L'$ and $R'$ are the least and largest entries of parity $p$ preceding $y'$.
\end{lemma}

\begin{proof}
If $2y>L+R$, then $z=2y-L>R$ does not precede $y$, so $w=3y-2L\in W$ by Lemma~\ref{lem:escape}; but $w$ has parity $1-p$ and $w>y$.
\end{proof}

\section{Computations}\label{sec:comp}

All computations below are corroborations and are not used in the proofs.

\begin{itemize}[leftmargin=*]
\item \emph{Theorem~\ref{thm:odd}.} With the enumerations $y=0,2,-2,4,-4,\dots$ and $t=1,-1,3,-3,\dots$, the word after $200$ steps has $40002$ terms, all values in $[-39601,40001]$, and it contains every integer of $[-199,199]$. An exhaustive check over all pairs of positions found no in-order $3$-AP of the form (odd, even, odd) and no monotone $4$-AP of odd difference. It found $261{,}437{,}641$ in-order $3$-APs, namely $261{,}436{,}991$ of even difference and $650$ of odd difference (all of the form (even, odd, even)), and $154{,}981{,}918$ monotone $4$-APs, all of even difference. This illustrates that Theorem~\ref{thm:odd} says nothing about even differences.
\item \emph{Safe words.} Safe words containing $[-n,n]$ exist for every $n$: by (T1) and (T2), a finite set listed in decreasing $\prec$-order is $\prec$-safe for every tree order $\prec$ (this is the argument of \cite[Lemma~1]{Ho26}, where $\prec\,=\,\lhd$ and the setting is $\N$). The question is extendability. We encoded the statement ``$W'$ extends $W$, contains a given set, uses only integers of a given window, and is $\prec$-safe for some tree order $\prec$'' as a SAT problem, with the tree order among the unknowns and all integers outside the window in the tail. The word
\[
7,-1,3,-5,5,-3,1,-7,6,-2,2,-6,4,-4,0,8,-8
\]
is the decreasing listing of $[-8,8]$ for the tree order that prefers digit $0$ at the nodes of levels $0,1,2$ and digit $1$ at the nodes of levels $3,4$. It has no safe extension that contains $[-9,9]$ and uses only integers in $[-64,64]$. Likewise, Ho's word $0,1,3,2,11,7,15,9,5,6,4$ for $\N$ has no extension containing $\{0,\dots,8\}$ inside $[0,96]$ that is safe for a tree order (counting only progressions in $\N$), whereas the next step of Ho's recursion (in our implementation) adds $8$ and produces a word of $354$ terms with largest entry $2623$. So bounded searches of this kind cannot decide extendability. Both unsatisfiability results were confirmed with a second, independently written encoding.
\end{itemize}

\section{Concluding remarks}

Problem~195 remains open. The results above show what a $2$-adic construction for $\Z$ must avoid. Protection at a node and at its unprotected child, or at a node and the grandchildren below its unprotected child, is impossible (Theorem~B). Safety cannot hold at all prefixes (Theorem~C). And the parity-block checkpoints of Ho's recursion collapse to midpoints (Theorem~D). Lemma~\ref{lem:largest} leaves room for constructions in which elements of the non-preferred parity are placed far outside the current range \emph{before} the preferred elements for which they serve as escapes. We do not know whether such ``anchored'' constructions exist.

\begin{question}
Is there a permutation of $\Z$ without monotone $4$-APs whose difference is not divisible by $4$?
\end{question}

\begin{question}
For which sets $K\subseteq\N$ is there a permutation of $\Z$ without monotone $4$-APs whose difference $d$ satisfies $v_2(d)\in K$? By Corollary~\ref{cor:level} every singleton works. By Theorems~\ref{thm:twolevel} and~\ref{thm:level02} and their renormalised versions, symmetric insertion at all nodes of the levels in $K$ fails as soon as $K$ contains two integers that differ by $1$ or by $2$. We do not know what happens for $K=\{0,3\}$.
\end{question}

\begin{thebibliography}{99}
\bibitem{Ad22} S.~Adenwalla, Avoiding monotone arithmetic progressions in permutations of integers, \emph{Discrete Math.} (2024), 114183; arXiv:2211.04451.
\bibitem{Ad23} S.~Adenwalla, A generalisation of a result on monotone arithmetic progressions in permutations of the positive integers, arXiv:2302.09662 (2023).
\bibitem{ABJ} H.~Ardal, T.~Brown, V.~Jungi\'c, Chaotic orderings of the rationals and reals, \emph{Amer. Math. Monthly} \textbf{118} (2011), 921--925.
\bibitem{DEGS} J.~A.~Davis, R.~C.~Entringer, R.~L.~Graham, G.~J.~Simmons, On permutations containing no long arithmetic progressions, \emph{Acta Arith.} \textbf{34} (1977), 81--90.
\bibitem{EP195} T.~F.~Bloom, Erd\H{o}s Problem \#195, \url{https://www.erdosproblems.com/195}.
\bibitem{FC195} Google DeepMind, Formal Conjectures, \texttt{ErdosProblems/195.lean}.
\bibitem{Ge19} J.~Geneson, Forbidden arithmetic progressions in permutations of subsets of the integers, \emph{Discrete Math.} \textbf{342} (2019), 1489--1491.
\bibitem{Ge26} J.~Geneson, Density bounds for permutations avoiding monotone arithmetic progressions, arXiv:2608.12604 (2026).
\bibitem{Ho26} B.~S.~Ho, A 4AP-free permutation of the positive integers, arXiv:2609.12780 (2026).
\bibitem{LV11} T.~D.~LeSaulnier, S.~Vijay, On permutations avoiding arithmetic progressions, \emph{Discrete Math.} \textbf{311} (2011), 205--207; arXiv:1004.1740.
\end{thebibliography}

\end{document}
